% Vietnamese translation for OI Public Library
% Source-File: IOI中国国家候选队论文/2003/刘才良/刘才良.ppt
\documentclass[11pt]{article}
\usepackage{oipl-vn}

\title{Bản trình chiếu: Ghi chú về đồ thị phẳng và ứng dụng}
\author{Liu Cailiang}
\date{2003}

\begin{document}
\maketitle

\origtitle{Slide companion for planar graph techniques}
\sourcepath{IOI China candidate team papers / 2003 / Liu Cailiang / Liu Cailiang.ppt}

\section{Phạm vi bản dịch}

Tệp này chỉ có bản PowerPoint. Văn bản khôi phục được gồm nhiều công thức đồ
thị phẳng, tên bài \textit{CEPC 2001}, \textit{CEOI 1997}, và một bài ACM
Shanghai 2000. Vì phần chữ Hán trong slide không đọc được bằng công cụ hiện
có, bản dịch ghi lại các ý thuật toán chắc chắn nhận ra từ công thức.

\section{Công thức nền}

Với đồ thị phẳng \(G=\langle V,E,F\rangle\), các slide dùng công thức Euler
\[
  |V|-|E|+|F|=2.
\]
Từ đó suy ra các cận quen thuộc như \(|E|\le 3|V|-6\) trong trường hợp mặt có
bậc ít nhất ba, và \(|E|\le 2|V|-4\) trong trường hợp đồ thị hai phía phẳng.

\section{Ví dụ ứng dụng}

Các slide nhắc tới bài \textit{CEPC 2001} với \(N\le 8000\), trong đó cần tổ
chức các đoạn theo tọa độ \(y\) và duy trì các danh sách \(C[2y]\),
\(C[2y+1]\). Một ví dụ khác là \textit{CEOI 1997} với \(N\le 500\), nơi các
đường đi kiểu \(1-3-5-6-12-10\) và \(9-11\) được dùng để minh họa cấu trúc
phẳng.

\section{Thông điệp thuật toán}

Khi một bài toán có mô hình phẳng, công thức Euler biến các đại lượng tưởng
như lớn thành tuyến tính theo \(|V|\). Điều này cho phép thay các thuật toán
\(O(N^3)\) hoặc \(O(N^2)\) bằng các cách làm gần tuyến tính hay
\(O(N\log N)\), tùy bước sắp xếp và cấu trúc dữ liệu đi kèm.

\end{document}
